\par\medskip
现转向一般情形. 考察 $X$ 的定义方程, 可设 $k$ 在 $\mathbb F_p$ 上有限生成.
于是存在子环 $B\subset k$, 它作为 $\mathbb F_p$ 代数有限生成且分式域为 $k$, 使得
\begin{equation}\tag{1.3.9}
X\subset\mathbb P_k^{n+r}
\end{equation}
是相对完全交 $\widetilde X\subset\mathbb P_B^{n+r}$ 的泛纤维,
且其每个纤维与 $X$ 多重次数相同并几何不可约.

记
\begin{equation}\tag{1.3.10}
\rho:\widetilde X\longrightarrow S=\operatorname{Spec}(B)
\end{equation}
为投影,
\begin{equation}\tag{1.3.11}
E=\rho_*\mathcal O_{\widetilde X}(d)
\end{equation}
为局部自由 $\mathcal O_S$ 模, 并记
\begin{equation}\tag{1.3.12}
\pi:T=\check{\mathbb P}(E)\longrightarrow S
\end{equation}
为相应的 $E$ 中直线丛. (1.3.12) 的泛纤维恰为
\[
U=\mathbb P\bigl(H^0(X,\mathcal O_X(d))\bigr).
\]
在 $T$ 上有 $\widetilde X\subset\mathbb P_B^{n+r}$ 的 $d$ 次超曲面泛截面, 记为
\begin{equation}\tag{1.3.13}
\alpha:\widetilde X.H_T\longrightarrow T.
\end{equation}
这些构造给出下列由笛卡儿方块组成的图:
\begin{equation}\tag{1.3.14}
\begin{tikzcd}[column sep=huge,row sep=large]
\widetilde X.H_T \arrow[d,"\alpha"]
  & X.H_u=(\widetilde X.H_T)\times_T U \arrow[l]\arrow[d] \\
T \arrow[d,"\pi"]
  & U=T\times_S\operatorname{Sp}(k) \arrow[l]\arrow[d] \\
S & \operatorname{Sp}(k)\arrow[l]
\end{tikzcd}
\end{equation}
$\mathcal O_T$ 模
\begin{equation}\tag{1.3.15}
R^{n-1}\alpha_*\mathcal O_{\widetilde X.H_T}
\end{equation}
在由满足 $\dim\alpha^{-1}(t)=n-1$ 的 $t$ 组成的开子集 $T'\subset T$ 上局部自由,
且在此处其形成与任意基变换可交换. 由于 $T$ 和 $U$ 具有同一泛点,
要证明 (1.3.4) 非空, 只需证明 $T'$ 的开子集
\begin{equation}\tag{1.3.16}
\bigl\{t\in T'\mid\widetilde X_{\pi(t)}.H_t
   \text{ 非特殊}\bigr\}
\end{equation}
非空. 但已证的情形 (1.3.5) 表明, 对每个闭点 $s\in S$, 纤维 $\pi^{-1}(s)$
与开集 (1.3.16) 的交非空. 定理得证.

\medskip
\noindent\emph{底本注.}
照录法文层保留了印本右上角的对象 $(\widetilde X.H_T)\times_U T$
及措辞 \emph{se formation}. 笛卡儿方块、$T$ 的定义及其泛纤维 $U$ 的定义
要求采用上文的 $(\widetilde X.H_T)\times_T U$ 和 \emph{sa formation}.
普通字形的 $E$、$H^0$ 中的指数 $0$ 以及定义 $U$ 的外括号,
都是直接根据扫描件所作的转录修正, 而非编辑校改.
